% Two-column preprint. Every figure, table and procedure is one column wide.
\documentclass[10pt,twocolumn,a4paper]{article}
\usepackage[margin=0.75in]{geometry}
\usepackage[T1]{fontenc}
\usepackage{lmodern}
% Shared definitions, adapted from the StaleyNET paper.

\usepackage{amsmath,amssymb,amsthm,mathtools}
\usepackage{graphicx}
\usepackage{booktabs}
\usepackage{array}
\usepackage[table]{xcolor}
\usepackage{tikz}
\usetikzlibrary{arrows.meta,positioning,calc,fit,shadows.blur,backgrounds}
\usepackage[most]{tcolorbox}
\usepackage{enumitem}
\usepackage{microtype}
\usepackage{xspace}

\definecolor{accent}{HTML}{3B4BA8}
\definecolor{definitionfill}{HTML}{EEF0FB}
\definecolor{codefill}{HTML}{F6F7F9}
\definecolor{rulefill}{HTML}{E7F5EF}
\definecolor{muted}{HTML}{5F6475}

% Syntax colors for pseudocode.
\definecolor{synkeyword}{HTML}{3B4BA8}
\definecolor{synstring}{HTML}{2E7D4F}

% Material tonal containers for figures: a fill and the dark text color of the same hue.
\definecolor{bluefill}{HTML}{D6E3FF}   \definecolor{blueink}{HTML}{001B3E}
\definecolor{violetfill}{HTML}{EADDFF} \definecolor{violetink}{HTML}{21005D}
\definecolor{greenfill}{HTML}{C3EFCB}  \definecolor{greenink}{HTML}{00210B}
\definecolor{amberfill}{HTML}{FFDDB3}  \definecolor{amberink}{HTML}{2A1700}
\definecolor{redfill}{HTML}{FFDAD6}    \definecolor{redink}{HTML}{410002}
\definecolor{grayfill}{HTML}{E9EAF0}   \definecolor{grayink}{HTML}{1B1C21}
\definecolor{edgeink}{HTML}{44474E}

% Definitions in shaded boxes.
\theoremstyle{definition}
\newtheorem{definition}{Definition}[section]
\tcolorboxenvironment{definition}{colback=definitionfill,colframe=definitionfill,boxrule=0pt,arc=2pt,left=6pt,right=6pt,top=4pt,bottom=4pt}

% The example line inside a definition.
\newcommand{\example}{\par\smallskip\noindent\textit{Example.}\ }

% A highlighted observation.
\newtcolorbox{observation}{colback=rulefill,colframe=rulefill,boxrule=0pt,arc=2pt,left=6pt,right=6pt,top=4pt,bottom=4pt,
  title={\sffamily\bfseries\footnotesize KEY OBSERVATION},coltitle=accent,colbacktitle=rulefill,
  attach title to upper,after title={\par\smallskip}}

% Pseudocode in the style of Cormen, Leiserson, Rivest and Stein: one
% statement per numbered line, procedure names in small capitals, and a
% wider gap after the last line of each group (end such a line with \groupend).
\newcounter{codeline}
\ExplSyntaxOn
\seq_new:N \l__proc_words_seq
\NewDocumentCommand{\proc}{m}{
  \textsf{\bfseries
    \seq_set_split:Nnn \l__proc_words_seq { - } { #1 }
    \seq_map_indexed_inline:Nn \l__proc_words_seq
      {
        \int_compare:nT { ##1 > 1 } { - }
        \tl_head:n { ##2 }
        \scalebox{0.8}{ \text_uppercase:n { \tl_tail:n { ##2 } } }
      }
  }
}
\ExplSyntaxOff

% A procedure in a shaded box. Arguments: the signature, the input in words, the output in
% words, and the link to the Rust lines that implement it. A procedure is one column wide.
\newcounter{procnum}
\newenvironment{procedure}[4]{%
  \setcounter{codeline}{0}\refstepcounter{procnum}%
  \begin{tcolorbox}[enhanced,colback=codefill,colframe=codefill,boxrule=0pt,arc=4pt,left=8pt,right=6pt,top=5pt,bottom=4pt]%
  \small\linespread{1.15}\selectfont\def\procsignature{#1}%
  \noindent{\sffamily\footnotesize\color{accent}Procedure \theprocnum}\hfill{\footnotesize #4}\par\vspace{2pt}%
  {\footnotesize\color{muted}\noindent\textit{Input:} #2\par\noindent\textit{Output:} #3\par}\vspace{5pt}%
  \begin{tabbing}%
  \hspace*{1.5em}\=\hspace{1.2em}\=\hspace{1.2em}\=\hspace{1.2em}\=\hspace{1.2em}\=\hspace{1.2em}\=\kill
  \hspace*{1.5em}\procsignature\ \{\\[1pt]}{%
  \\[1pt]\hspace*{1.5em}\}%
  \end{tabbing}\end{tcolorbox}}
\newcommand{\li}{\stepcounter{codeline}\llap{\scriptsize\color{muted}\thecodeline\hspace{0.6em}}\>\>}
\newcommand{\groupend}{\\[0.6\baselineskip]}
\newcommand{\kw}[1]{{\color{synkeyword}\textbf{#1}}}
\newcommand{\cmt}[1]{\hfill{\color{muted}\itshape\footnotesize // #1}}
\newcommand{\str}[1]{{\color{synstring}\texttt{"#1"}}}
\newcommand{\var}[1]{\ensuremath{\mathit{#1}}}
\newcommand{\st}[1]{\textsf{#1}}
\newcommand{\nil}{\textsc{nil}}
% A Rust type in a procedure.
\definecolor{syntype}{HTML}{7B3FA0}
\newcommand{\ty}[1]{{\color{syntype}\textsf{#1}}}

% Links to exact lines of the harness in the public repository.
\newcommand{\repo}{https://github.com/arpanpathak/anti_ai_slop}
\newcommand{\commit}{main}

\newcommand{\srcref}[3]{\href{\repo/blob/\commit/harness/src/#1\#L#2-L#3}{\texttt{#1}:#2--#3}}

% Rust listings with syntax colors.
\usepackage{listings}
\definecolor{synmacro}{HTML}{A1520B}
\lstdefinestyle{rust}{
  basicstyle=\fontsize{6.5}{7.8}\selectfont\ttfamily,
  keywords={fn,let,mut,match,loop,enum,impl,struct,self,Self,return,if,pub,const,use,mod},
  keywordstyle=\color{synkeyword}\bfseries,
  keywords=[2]{State,Result,Error,Output,Path,Command,OsStr,Answer,Harness,String,Ok,Err,Some},
  keywordstyle=[2]\color{syntype},
  keywords=[3]{bail,ensure,format,println,eprintln,json},
  keywordstyle=[3]\color{synmacro},
  comment=[l]{//},
  commentstyle=\color{muted}\itshape,
  morestring=[b]",
  stringstyle=\color{synstring},
  showstringspaces=false,
  columns=fullflexible,
  keepspaces=true,
  numbers=left,
  numberstyle=\tiny\color{muted},
  numbersep=5pt,
  xleftmargin=10pt,
  breaklines=true,
  breakindent=1.5em,
  postbreak=\mbox{{\color{muted}$\hookrightarrow$}\space}
}

% A one-column listing copied from a source file, with a link to the same lines on GitHub.
% Arguments: file in harness/src/, first line, last line, label, one-line title.
\newcounter{listnum}
\newcommand{\listingcard}[5]{%
  \begin{procfloat}[!t]%
  \begin{tcolorbox}[enhanced,colback=codefill,colframe=codefill,boxrule=0pt,arc=4pt,left=6pt,right=6pt,top=5pt,bottom=3pt]%
  \refstepcounter{listnum}\label{#4}%
  {\sffamily\footnotesize\color{accent}Listing \thelistnum}\hfill{\footnotesize\srcref{#1}{#2}{#3}}\par
  {\small #5}\par\vspace{2pt}%
  \lstinputlisting[style=rust,firstline=#2,lastline=#3,firstnumber=1]{../../harness/src/#1}%
  \end{tcolorbox}%
  \end{procfloat}}

% Procedures float in one column, like figures.
\makeatletter
\newcommand\fps@procfloat{tbp}
\newcommand\ftype@procfloat{4}
\newcommand\ext@procfloat{lop}
\newenvironment{procfloat}{\@float{procfloat}}{\end@float}
\makeatother

% Figures: Material cards with tonal fills, no outlines, a soft shadow below each box, and
% sharp arrowheads. The shadow is offset further down than it is blurred, so it never shows
% above a box as a ring.
\tikzset{
  figure/.style={font=\sffamily\small, >={Stealth[length=2.5mm, width=1.9mm]}, line width=0.75pt},
  elevated/.style={blur shadow={shadow xshift=0pt, shadow yshift=-1.6pt,
    shadow blur radius=1.4pt, shadow blur steps=14, shadow opacity=40,
    shadow blur extra rounding=1.4pt}},
  block/.style={rounded corners=6pt, draw=none, align=center, inner sep=5pt,
    minimum height=2.3em, elevated},
  chip/.style={rounded corners=3pt, draw=none, align=center, inner sep=2pt,
    font=\sffamily\scriptsize, minimum height=1.55em},
  region/.style={rounded corners=10pt, inner sep=8pt},
  region label/.style={font=\sffamily\scriptsize\bfseries, text=edgeink},
  edge label/.style={font=\sffamily\footnotesize, text=edgeink, align=center},
  flow/.style={->, draw=edgeink, line cap=butt},
}

% A figure card: a title and a one-line lede inside the figure, then the drawing.
\newcommand{\figcard}[3]{%
  \begin{tcolorbox}[enhanced,colback=white,colframe=white,boxrule=0pt,arc=6pt,left=7pt,right=7pt,top=6pt,bottom=7pt,
    drop fuzzy midday shadow=black!30]%
  {\sffamily\bfseries\small #1}\par\vspace{1pt}{\sffamily\footnotesize\color{muted}#2}\par\vspace{7pt}%
  \centering #3%
  \end{tcolorbox}}

\newcommand{\code}[1]{\texttt{#1}}

% A named constant in a procedure and in the text.
\newcommand{\const}[1]{\textsf{#1}}

\usepackage[numbers,sort&compress]{natbib}
\usepackage[colorlinks,linkcolor=accent,citecolor=accent,urlcolor=accent]{hyperref}
\usepackage[capitalise,noabbrev]{cleveref}
\crefname{procnum}{Procedure}{Procedures}
\crefname{listnum}{Listing}{Listings}
\setlength{\bibsep}{0pt}
\renewcommand{\bibfont}{\footnotesize}
\setlength{\columnsep}{0.25in}
\setlength{\emergencystretch}{3em}
\hyphenation{sto-chas-tic de-ter-min-is-tic name-space name-spaces}
\renewcommand{\topfraction}{0.9}
\renewcommand{\textfraction}{0.08}
\renewcommand{\floatpagefraction}{0.85}
\raggedbottom

\title{\textbf{A State-Machine Harness and a Namespace Sandbox\\for Code Written by Language Models}}
\author{Arpan Pathak}
\date{November 2026}

\begin{document}
\maketitle

\begin{abstract}
When coding agents execute programs written by language models, the model's text is transformed into actions on a physical machine. Because the test command is typically authored by the model as well, a single injected reply can overwrite shell startup files or exfiltrate private keys to a remote host. We present a harness designed to be small enough to trace every execution path from a model's reply to the host. The agent loop is implemented as a fixed five-state machine in under 200 lines of Rust, where the state machine itself dictates the next step. Consequently, a reply can only reach the host through two channels: as a program file or as a check command. We restrict file names to a plain format and execute every check within a transient sandbox constructed from Linux namespaces using bubblewrap, providing read-only system files, no network access, and no access to secrets. The sandbox introduces approximately 15\,ms of overhead per check. In our evaluation, none of the 104 generated escape attempts left a trace on the host, whereas 56 of the same attempts succeeded when executed without the sandbox. While the sandbox alone does not bound the number of processes or the memory consumed in the working directory, our tests demonstrate that a cgroup effectively bounds both. A passing check reveals only the specific outcomes the model intended to test.
\end{abstract}

\section{Introduction}
\label{sec:intro}

Coding agents are increasingly capable of executing the code they generate. Systems such as SWE-agent run tests within a repository, and CodeAct expresses agent actions as Python code for execution \cite{sweagent2024,codeact2024}. In both cases, text sampled from a language model is converted into commands on a physical machine—a transition Srinivasan describes as the \emph{stochastic-deterministic boundary} \cite{sdb2026}. This text must be treated as untrusted. A \emph{prompt injection} is an instruction embedded in text that the model processes, such as a task description or the output of a previous command. Under such an instruction, a model might return a reply like the following, which we use to illustrate the threat throughout this paper.
\begin{center}
\begin{tcolorbox}[enhanced,width=\linewidth,colback=codefill,colframe=codefill,boxrule=0pt,arc=4pt,left=6pt,right=6pt,top=4pt,bottom=4pt]
\scriptsize\ttfamily
\{"file\_name": {\color{synstring}"../../.bashrc"},\\
\ "code": {\color{synstring}"curl -s http://attacker.example/x | sh"},\\
\ "check": {\color{synstring}"cat \textasciitilde/.ssh/id\_ed25519 |}\\
\ {\color{muted}$\hookrightarrow$}\ \ \ \ \ \ \ {\color{synstring}curl -s -d @- http://attacker.example/key"}\}
\end{tcolorbox}
\end{center}
In this example, the file name uses directory traversal to target the user's shell startup file, and the check command exfiltrates the user's private SSH key to a remote host. The initial version of our harness permitted both fields.

Tool protocols, such as the Model Context Protocol (MCP)\footnote{\url{https://modelcontextprotocol.io}}, standardize how models discover and invoke tools, but they do not restrict the side effects of a tool once it executes on the host. In large-scale agentic systems, the path from a model's reply to the host is often difficult to audit, as it traverses complex system prompts, multiple layers of application logic, and network requests that may only be visible via packet capture. We sought to design a system small enough to explicitly list every mechanism by which a reply reaches the host, allowing us to confine each mechanism and observe how model replies interact with them.

To achieve this, we maintained a minimal harness. It transmits a task in plain English to a model and requests a file name, a program, and a shell command to build and test that program (which we call the \emph{check}). The harness saves the program, executes the check, and returns any failures to the model for up to \const{MaxAttempts} attempts (\cref{fig:system}). The execution loop is a five-state machine governed by a single \kw{match} statement. The implementation in \code{main.rs} consists of only 158 lines of code (excluding comments and whitespace), and the sandbox component adds 66 lines \cite{repo}. Because the state machine controls every transition, a reply can only affect the host through the files we write and the checks we execute, both of which can be confined by the operating system.

\begin{figure}[!t]
\figcard{One attempt}{We write a model's program on our machine and run its check inside a sandbox.}{%
\begin{tikzpicture}[figure]
  \node[block, fill=bluefill, text=blueink, minimum width=3.6cm] (model) at (0,0)
    {\textbf{model server}\\[-1pt]\scriptsize chat API};
  \node[block, fill=violetfill, text=violetink, minimum width=3.6cm] (harness) at (0,-2.1)
    {\textbf{harness}\\[-1pt]\scriptsize state machine of \cref{fig:states}};
  \node[block, fill=amberfill, text=amberink, minimum width=3.0cm] (program) at (-1.95,-4.6)
    {\texttt{\scriptsize harness\_out/main.rs}\\[-1pt]\scriptsize plain file name};
  \node[block, fill=white, text=grayink, minimum width=3.0cm] (check) at (1.95,-4.6)
    {\texttt{\scriptsize rustc main.rs}\\[-2pt]\texttt{\scriptsize \&\& ./main}};
  \coordinate (sandbox top) at ($(check.north)+(0,0.45)$);
  \begin{scope}[on background layer]
    \node[region, fill=greenfill!80, fit=(check)(sandbox top), inner sep=5pt] (sandbox) {};
  \end{scope}
  \node[region label, text=greenink, anchor=north east] at ($(sandbox.north east)+(-2mm,-1.2mm)$) {sandbox};
  \coordinate (send) at ($(harness.north)+(-0.7,0)$);
  \coordinate (receive) at ($(harness.north)+(0.7,0)$);
  \coordinate (write) at ($(harness.south)+(-1.2,0)$);
  \coordinate (result) at ($(check.north)+(-0.75,0)$);
  \draw[flow] (send) -- node[edge label, left] {conversation} (send |- model.south);
  \draw[flow] (receive |- model.south) -- node[edge label, right] {reply} (receive);
  \draw[flow] (write) -- node[edge label, left] {write} (write |- program.north);
  \draw[flow] (program) -- node[edge label, above] {bind} (check);
  \draw[flow] (result) -- node[edge label, right, pos=0.6] {exit status\\[-2pt]and output} (result |- harness.south);
\end{tikzpicture}}
\caption{One attempt, with a check from our sample run in \cref{sec:evaluation}. Our model server runs on a separate machine.}
\label{fig:system}
\end{figure}


Our procedures follow the conventions of Cormen et al.\ \cite{cormen2022introduction}, utilizing Rust types for variables and providing direct links to their implementations. Our contributions are:
\begin{itemize}[nosep,leftmargin=*]
  \item a harness featuring an agent loop implemented as a fixed state machine, limited to \const{MaxAttempts} requests per task (\cref{sec:harness});
  \item a threat model that traces each field of a reply to its potential impact on the host (\cref{sec:threats});
  \item a sandbox for model-generated checks, constructed from Linux namespaces without requiring root privileges, featuring a file name restriction and a startup probe (\cref{sec:sandbox});
  \item an evaluation comprising 104 generated escape attempts, a test of cgroup limits, an analysis of sandbox overhead, and a real-world task (\cref{sec:evaluation}).
\end{itemize}

\section{The harness}
\label{sec:harness}

Each conversation begins with instructions that enforce a specific reply format. The instructions also require the program to assert its own result, ensuring that an incorrect answer causes its own check to fail.

\begin{definition}[Reply]\label{def:reply}
A \emph{reply} is a JSON object containing the following string fields: \code{file\_name}, \code{code}, and \code{check}. These represent the name of the program file, the complete program source, and a single shell command used to build and execute the program, respectively.
\example
{\small\ttfamily
\{"file\_name": {\color{synstring}"main.rs"},\\
\ "code": {\color{synstring}"fn main() \{ ... \}"},\\
\ "check": {\color{synstring}"rustc main.rs \&\& ./main"}\}}
\end{definition}

\begin{procfloat}[!t]
\begin{procedure}{\proc{Harness}$(\var{task}\colon$ \ty{String}$)\ \rightarrow$ \ty{Result}\label{proc:harness}}
{a task in plain English.}
{\code{Ok} with the reply whose check passed, or \code{Err} after \const{MaxAttempts} failed attempts.\par\noindent
\textit{Types:} \ty{State} $=$ \st{Ask} $\mid$ \st{Check}(\ty{String}) $\mid$ \st{Fix}(\ty{Error}) $\mid$ \st{Pass} $\mid$ \st{GiveUp}.}
{\srcref{main.rs}{64}{151}}
\li \kw{if} \proc{Run-Sandboxed}$(\str{true}, \nil)$ \kw{is} \code{Err}$(\var{error})$ \{ \\
\li \> \kw{return} \code{Err}$(\var{error})$ \cmt{no sandbox on this host} \\
\li \} \groupend
\li $\var{conversation}\colon$ \ty{Conversation} $= \langle \const{Instructions}, \var{task} \rangle$ \\
\li $\var{attempt}\colon$ \ty{u32} $= 0$ \\
\li $\var{state}\colon$ \ty{State} $=$ \st{Ask} \groupend
\li \kw{loop} \{ \\
\li \> \kw{match} $\var{state}$ \{ \\
\li \> \> \st{Ask} $\Rightarrow$ \{ \\
\li \> \> \> $\var{attempt} = \var{attempt} + 1$ \\
\li \> \> \> $\var{state} = \proc{Ask}(\var{conversation})$ \\
\li \> \> \} \groupend
\li \> \> \st{Check}$(\var{reply})$ $\Rightarrow$ \{ \\
\li \> \> \> $\var{state} = \proc{Sandboxed-Check}(\var{reply})$ \\
\li \> \> \} \groupend
\li \> \> \st{Fix}$(\var{failure})$ $\Rightarrow$ \{ \\
\li \> \> \> append $\var{failure}$ to $\var{conversation}$ \\
\li \> \> \> $\var{state} =$ \st{Ask} \groupend
\li \> \> \> \kw{if} $\var{attempt} = \const{MaxAttempts}$ \{ \\
\li \> \> \> \> $\var{state} =$ \st{GiveUp} \\
\li \> \> \> \} \\
\li \> \> \} \groupend
\li \> \> \st{Pass} $\Rightarrow$ \kw{return} \code{Ok}$(\var{reply})$ \groupend
\li \> \> \st{GiveUp} $\Rightarrow$ \kw{return} \code{Err}$(\str{no passing program})$ \\
\li \> \} \\
\li \}
\end{procedure}
\end{procfloat}


\Cref{proc:harness} describes the execution loop, and \cref{fig:states} illustrates it as a state machine. The model does not determine the next state; instead, a reply is treated as data that a single \kw{match} statement inspects to decide the next transition. The only transition dependent on the model's output leaves \st{Check} and moves toward either \st{Pass} or \st{Fix}; all other transitions are fixed. In Rust (\cref{lst:run}), \ty{State} is an enum with five variants, and the compiler ensures that the \kw{match} statement handles every possible state. Consequently, adding a new state requires explicit handling in every part of the loop.

\begin{figure}[!t]
\figcard{States of the harness}{Only Ask sends a request to a model. Each box names the lines of \cref{proc:harness} that it runs.}{%
\begin{tikzpicture}[figure, state/.style={block, minimum width=1.6cm, minimum height=1.0cm}]
  \node[state, fill=bluefill, text=blueink] (ask) at (0,0) {\textbf{Ask}\\[-2pt]\scriptsize lines 10--11};
  \node[state, fill=violetfill, text=violetink] (check) at (2.9,0) {\textbf{Check}\\[-2pt]\scriptsize line 14};
  \node[state, fill=greenfill, text=greenink] (pass) at (5.8,0) {\textbf{Pass}\\[-2pt]\scriptsize line 23};
  \node[state, fill=amberfill, text=amberink] (fix) at (2.9,-2.1) {\textbf{Fix}\\[-2pt]\scriptsize lines 17--21};
  \node[state, fill=redfill, text=redink] (giveup) at (5.8,-2.1) {\textbf{GiveUp}\\[-2pt]\scriptsize line 24};
  \draw[flow] (ask) -- node[edge label, above] {reply} (check);
  \draw[flow] (check) -- node[edge label, above] {passed} (pass);
  \draw[flow] (check) -- node[edge label, right] {failed} (fix);
  \draw[flow] (fix) -- node[edge label, above] {last\\[-2pt]attempt} (giveup);
  \draw[flow] (fix.west) -| node[edge label, below, pos=0.3] {attempts\\[-2pt]left} (ask.south);
\end{tikzpicture}}
\caption{The 5 states of \cref{proc:harness} and their transitions.}
\label{fig:states}
\end{figure}


\listingcard{main.rs}{93}{105}{lst:run}{The implementation of \cref{proc:harness} in Rust. Each state's method returns the subsequent state.}

The loop is also bounded in duration. The single request to the model is made on line 11, and the attempt counter is incremented on line 10. The \st{Fix} state returns to \st{Ask} only if the counter remains below \const{MaxAttempts} (lines 19--21). Thus, a single task consumes at most \const{MaxAttempts} requests and an equivalent number of checks, regardless of the model's output.

Before the initial request, lines 1--3 execute an empty check (\code{true}) within the sandbox. If the sandbox fails to initialize on a given host, the harness terminates immediately. This probe prevents the model from wasting its allotted attempts on errors caused by a malfunctioning sandbox rather than its own code.

\section{Threat model}
\label{sec:threats}

We assume an attacker controls the entire content of every reply and can modify it during any attempt—for instance, via an injection in the task description or the output of a previous check. We assume the harness, the kernel, and the toolchain are trustworthy.

\begin{observation}
Because the state machine dictates every step, a reply can reach the host through exactly two channels: as the file written by the harness and as the check command it executes. The program itself runs exclusively within its check.
\end{observation}

Our initial implementation left both channels unprotected. It concatenated the file name to the working directory using a standard path join, which preserves \code{..} components. When running in \code{/home/user/harness}, this allowed writing \code{../../.bashrc} to \code{/home/user/.bashrc}. Furthermore, the check was executed via \code{sh -c} with the user's own permissions, enabling commands like \code{cat} to read private keys and \code{curl} to exfiltrate them.

\begin{table}[!t]
\centering
\caption{Ways in which a reply can act on its host, and our defense against each one.}
\label{tab:threats}
\footnotesize
\renewcommand{\arraystretch}{1.12}
\begin{tabular}{@{}>{\raggedright\arraybackslash}p{3.85cm}>{\raggedright\arraybackslash}p{3.95cm}@{}}
\toprule
Threat & Defense \\
\midrule
A file name with a directory part, such as \code{../../.bashrc} & Refused before any write. \\
A check that reads the home folder or the API key & Neither exists in the sandbox. \\
A check that sends data over the network & Only a loopback interface exists. \\
A check that changes host files or leaves a link for a later write & A read-only system, and a working folder that disappears with each check. \\
A check that runs forever or uses too much memory & Limits on time, memory, file size and open files (\cref{tab:limits}). \\
A check that fills memory with many files & Open in the harness. A cgroup stops it (\cref{tab:cgroup}). \\
A check that starts processes in a loop & Open in the harness. A cgroup stops it (\cref{tab:cgroup}). \\
A check that uses a kernel bug & Open. We load no seccomp filter. \\
Output that carries instructions to the model & Open. We return its output unchanged. \\
A check that tests nothing & Open. A model writes its own check. \\
\bottomrule
\end{tabular}
\end{table}


\Cref{tab:threats} summarizes each way a reply can interact with the host and our corresponding defense. The file name restriction secures the first row, while the sandbox secures the subsequent four. The remaining rows are addressed in \cref{sec:evaluation,sec:limits}.

\section{The sandbox}
\label{sec:sandbox}

\subsection{Plain file names}

The file name is the simpler of the two ways in, because the harness does the writing. We only
need to make sure that a name cannot point outside \const{WorkDir}.

\begin{definition}[Plain file name]\label{def:plain}
A file name is \emph{plain} if it is not empty, is neither \code{.} nor \code{..}, and contains
no slash. A plain name always names a file directly inside the working folder.
\example \code{main.rs} is plain. \code{../../.bashrc} and \code{/etc/passwd} contain slashes.
\end{definition}

\begin{procfloat}[!t]
\begin{procedure}{\proc{Sandboxed-Check}$(\var{reply}\colon$ \ty{String}$)\ \rightarrow$ \ty{State}\label{proc:check}}
{the text of one reply.}
{\st{Pass}, or \st{Fix} with the failure that we send back to the model.\par\noindent
\textit{Types:} \ty{Answer} $= \{\,\var{file\_name}\colon$ \ty{String}, $\var{code}\colon$ \ty{String}, $\var{check}\colon$ \ty{String}$\,\}$.}
{\srcref{main.rs}{124}{225}}
\li $\var{answer}\colon$ \ty{Answer} $= \proc{Parse-Json}(\var{reply})$ \\
\li \kw{if} $\var{answer} = \nil$ \{ \\
\li \> \kw{return} \st{Fix}$(\str{not the JSON asked for})$ \\
\li \} \groupend
\li $\var{name}\colon$ \ty{String} $= \var{answer.file\_name}$ \\
\li \kw{if} $\var{name}$ is not plain \{ \\
\li \> \kw{return} \st{Fix}$(\str{not a plain file name})$ \\
\li \} \groupend
\li $\var{program}\colon$ \ty{Path} $=$ \const{WorkDir} $/\ \var{name}$ \\
\li write $\var{answer.code}$ to $\var{program}$ \groupend
\li $\var{check}\colon$ \ty{String} $= \var{answer.check}$ \\
\li $\var{output}\colon$ \ty{Output} $= \proc{Run-Sandboxed}(\var{check}, \var{program})$ \\
\li \kw{if} $\var{output.status} \neq 0$ \{ \\
\li \> \kw{return} \st{Fix}$(\var{output})$ \\
\li \} \groupend
\li \kw{return} \st{Pass}
\end{procedure}
\end{procfloat}


\Cref{proc:check} refuses a name that is not plain before it writes anything (lines 5 to 8). We
compare characters and do not resolve the path. Resolution follows symbolic links and depends on
what is on disk at that moment, while a test of characters gives the same answer every time. The
test is sufficient here because only the harness writes into \const{WorkDir} and no check can see
that folder, so no check can plant a link in it. On Linux, \code{openat2} with
\code{RESOLVE\_BENEATH} enforces a comparable rule inside the kernel \cite{openat2}.

\subsection{The view of a check}

The check is more difficult to confine because it is an arbitrary shell command. Since we cannot reliably predict the behavior of a shell command by static analysis, we instead restrict its visibility.

\begin{definition}[View of a check]\label{def:view}
$\proc{Run-Sandboxed}(\var{check}, \var{program})$ executes \var{check} using \code{sh} within a new sandbox containing only:
\begin{itemize}[nosep,leftmargin=*]
  \item read-only system directories and the toolchain;
  \item the program, read-only, at \code{/work/}\var{name};
  \item transient, in-memory \code{/work} and \code{/tmp} directories;
  \item a process table containing only the check's own processes;
  \item a network restricted to the loopback interface;
  \item the environment variables \code{PATH} and \code{HOME}${}={}$\code{/work}, plus necessary toolchain variables.
\end{itemize}
\example The check in \cref{def:reply} begins in \code{/work}, reads \code{/work/main.rs}, and writes the resulting binary to memory.
\end{definition}

\begin{figure}[!t]
\figcard{What a check can open}{Blue rows are read-only copies of host files. Red rows stay hidden, and green rows are new for each check.}{%
\begin{tikzpicture}[figure, row/.style={chip, minimum width=3.15cm, text width=3.0cm}]
  \node[region label] at (-1.95,0.52) {host};
  \node[region label] at (1.95,0.52) {sandbox};
  \node[row, fill=bluefill, text=blueink] (h1) at (-1.95,0) {\texttt{/usr}};
  \node[row, fill=bluefill, text=blueink] (h2) at (-1.95,-0.5) {\texttt{/home/user/.rustup}};
  \node[row, fill=bluefill, text=blueink] (h3) at (-1.95,-1.0) {\texttt{/home/user/.cargo/bin}};
  \node[row, fill=bluefill, text=blueink] (h4) at (-1.95,-1.5) {\texttt{harness\_out/main.rs}};
  \node[row, fill=redfill, text=redink] (h5) at (-1.95,-2.0) {\texttt{/home/user/.ssh}};
  \node[row, fill=redfill, text=redink] (h6) at (-1.95,-2.5) {\texttt{API\_KEY}, all variables};
  \node[row, fill=redfill, text=redink] (h7) at (-1.95,-3.0) {every host process};
  \node[row, fill=redfill, text=redink] (h8) at (-1.95,-3.5) {every network interface};
  \node[row, fill=bluefill, text=blueink] (s1) at (1.95,0) {\texttt{/usr}};
  \node[row, fill=bluefill, text=blueink] (s2) at (1.95,-0.5) {\texttt{/home/user/.rustup}};
  \node[row, fill=bluefill, text=blueink] (s3) at (1.95,-1.0) {\texttt{/home/user/.cargo/bin}};
  \node[row, fill=bluefill, text=blueink] (s4) at (1.95,-1.5) {\texttt{/work/main.rs}};
  \node[row, fill=greenfill, text=greenink] (s5) at (1.95,-2.0) {\texttt{/work}, \texttt{/tmp}: empty};
  \node[row, fill=greenfill, text=greenink] (s6) at (1.95,-2.5) {\texttt{PATH}, \texttt{HOME}, \texttt{RUSTUP\_HOME}};
  \node[row, fill=greenfill, text=greenink] (s7) at (1.95,-3.0) {processes of the check};
  \node[row, fill=greenfill, text=greenink] (s8) at (1.95,-3.5) {loopback interface \texttt{lo}};
  \foreach \i in {1,2,3,4} { \draw[flow] (h\i.east) -- (s\i.west); }
\end{tikzpicture}}
\caption{View of a check (\cref{def:view}), with our implementation's paths. Inside a sandbox, \texttt{HOME} is \texttt{/work}.}
\label{fig:mounts}
\end{figure}


\Cref{fig:mounts} illustrates this view. Each component of \cref{def:view} eliminates a potential attack vector. By clearing the home directory and environment, the check cannot access SSH or API keys. Because system directories are read-only and the \code{/work} directory is transient, no side effects from the check persist on the host. Mounting the program as read-only ensures that the file we preserve is exactly the one that passed the check. Finally, restricting the network to the loopback interface prevents connections to external hosts.

To verify this, we re-executed the hostile check using the plain name \code{main.sh}. Within the sandbox, \code{HOME} is set to \code{/work}, so \code{\textasciitilde/.ssh/id\_ed25519} resolves to \code{/work/.ssh/id\_ed25519}, which does not exist. The \code{cat} command reported \code{No such file or directory}, and \code{curl} exited with status 6, indicating a failure to resolve the host name.

The remainder of \cref{proc:check} manages standard execution failures. If a reply is not valid JSON according to \cref{def:reply} (e.g., if it is truncated by the model's output limit), the harness returns a short error message to the model (lines 1--4). If a check fails, the harness returns its exit status and combined output streams so that compiler errors are visible to the model (line 14).

\subsection{The sandbox on Linux}

A \emph{namespace} provides a group of processes with its own copy of a kernel resource, such as mounted file systems, process IDs, or network interfaces \cite{namespaces7}. Mount namespaces were introduced in Linux 2.4.19 in 2002. User namespaces, completed in Linux 3.8 in 2013, allow unprivileged users to create these environments \cite{kerrisk2013namespaces,userns7}. Bubblewrap, the sandboxing tool used by Flatpak, constructs views like \cref{def:view} using these namespaces \cite{bubblewrap}.

\begin{table}[!t]
\centering
\caption{Constants of our harness and limits of each check, with the values that we use.}
\label{tab:limits}
\footnotesize
\renewcommand{\arraystretch}{1.12}
\begin{tabular}{@{}l l >{\raggedright\arraybackslash}p{3.9cm}@{}}
\toprule
Name & Value & Meaning \\
\midrule
\const{MaxAttempts} & 5 & requests for one task \\
\const{WorkDir} & \code{harness\_out} & host folder for the programs \\
\const{WallSeconds} & 120\,s & wall-clock time of a check \\
\const{KillAfter} & 5\,s & from SIGTERM to SIGKILL \\
\const{CpuSeconds} & 60\,s & CPU time of each process \\
\const{AddressSpace} & 4\,GiB & memory of each process \\
\const{FileSize} & 64\,MiB & size of each file \\
\const{OpenFiles} & 256 & open files of each process \\
\bottomrule
\end{tabular}
\end{table}


\listingcard{sandbox.rs}{61}{86}{lst:command}{\proc{Run-Sandboxed} in Rust. The function collects the arguments of three programs and has no branch.}

We initiate each check using three standard programs, each with a single responsibility (\cref{lst:command}). GNU \code{timeout} sends \code{SIGTERM} after \const{WallSeconds} and \code{SIGKILL} \const{KillAfter} later. \code{prlimit} enforces the per-process limits specified in \cref{tab:limits}. Finally, \code{bwrap} creates new user, mount, PID, network, IPC, UTS, and cgroup namespaces to mount the view. It runs as process 1 of the new PID namespace; when it exits, the kernel terminates all processes started by the check. A new session prevents the check from interacting with our terminal.

We invoke bubblewrap as an external program rather than creating namespaces directly from Rust to avoid the complexities and pitfalls of system calls. For example, \code{unshare} with \code{CLONE\_NEWPID} does not move the caller into the new PID namespace, but only its subsequent children. A sandbox that failed to account for this would leave the check running in the host's PID namespace, where \code{kill -9 -1} could reach every process owned by the user.

Some hosts restrict unprivileged user namespaces. Ubuntu 24.04, for instance, blocks them by default via \path{kernel.apparmor_restrict_unprivileged_userns=1}. The probe at the start of \cref{proc:harness} detects this condition before the first request is made.

\section{Evaluation}
\label{sec:evaluation}

We conducted our measurements on an ARM64 board running Linux 5.15.148 and Ubuntu 22.04, using bubblewrap 0.6.1. The model, \code{Qwen3.6-35B-A3B}, was accessed via a chat API on a separate machine. Each experiment is implemented as a program in our repository, and every run is recorded with its corresponding commit and kernel version \cite{repo}.

\begin{table}[!t]
\centering
\caption{Escape attempts and observed results. Rows marked T are unit tests, and \code{/home/user} stands for our home folder.}
\label{tab:escapes}
\footnotesize
\renewcommand{\arraystretch}{1.12}
\begin{tabular}{@{}>{\raggedright\arraybackslash}p{3.55cm}>{\raggedright\arraybackslash}p{4.25cm}@{}}
\toprule
Check & Result \\
\midrule
T: file name \code{../escaped.sh} & Refused before any write. \\
T: write \code{left.txt} in \code{/work} & Written, and absent from \code{harness\_out} afterward. \\
T: overwrite the program & \code{Read-only file system}. \\
T: look for the project's \code{Cargo.toml} and \code{CARGO\_MANIFEST\_DIR} & Neither exists. \\
T: list the interfaces in \code{/proc/net/dev} & Only \code{lo}. \\
The hostile check of \cref{sec:intro} & \code{No such file or directory}, and \code{curl} exits with status 6. \\
\code{cat} \path{/home/user/.ssh/id_ed25519} & \code{No such file or directory}. \\
\code{curl http://1.1.1.1} & \code{Couldn't connect to server}. \\
\code{env} & Only \code{HOME}, \code{PATH}, \code{PWD} and \code{RUSTUP\_HOME}. \\
\code{ps -e}, then \code{kill -9 -1} & 3 processes: \code{bwrap}, \code{sh} and \code{ps}. The signal finds no process. \\
\code{dd} of 100\,MiB into \code{/work} & \code{File size limit exceeded} at 64\,MiB. \\
\code{sleep 300} & Ended by \code{timeout} after 120\,s, with status 124. \\
\bottomrule
\end{tabular}
\end{table}


\paragraph{Escape attempts.} We first manually authored 12 escape attempts, including 5 unit tests (\cref{tab:escapes}). None of these succeeded. While these 14 unit tests cover 93\% of the sandbox code, a manual list only explores single paths. To ensure a more robust evaluation, we generated variations. Each attempt targets one of the following: reading a secret file, creating a file, or connecting to a host listener. We tested four file-naming strategies: plain paths, \code{/proc/1/root}, \code{/proc/self/root}, and directory traversal via \code{..}. Using 6 programs for reading, 6 for writing, 2 directories, and 4 clients with 2 hostnames, we produced 104 distinct attempts.

We evaluated each attempt by observing its effect on the host. A read attempt was considered an escape if the secret appeared in the output; a write attempt succeeded if the target file existed; and a connection attempt was successful if our listener accepted it. To validate our detection method, we ran all 104 attempts in a standard shell without the sandbox. In that unconfined environment, 56 attempts succeeded. Inside the sandbox, none succeeded.

\begin{table}[!t]
\centering
\caption{Two checks that stay within every per-process limit, run in the sandbox alone and in
the same sandbox inside a cgroup of 64 tasks and 128\,MiB with no swap.}
\label{tab:cgroup}
\footnotesize
\renewcommand{\arraystretch}{1.12}
\begin{tabular}{@{}>{\raggedright\arraybackslash}p{2.5cm}>{\raggedright\arraybackslash}p{2.2cm}>{\raggedright\arraybackslash}p{2.75cm}@{}}
\toprule
Check & Sandbox & Sandbox in a cgroup \\
\midrule
Start 500 sleeping processes & 500 started & 59 started \\
Write 6 files of 60\,MiB to \code{/work} & 360\,MiB held in memory & Stopped by \code{SIGKILL} after 144\,ms \\
\bottomrule
\end{tabular}
\end{table}


\paragraph{Processes and memory.} The limits imposed by \code{prlimit} apply to individual processes, leaving two potential vulnerabilities: a check could spawn many processes, each within its own limits, or it could fill the \code{/work} directory with many small files that collectively consume significant memory. We tested both scenarios using the sandbox alone and using the sandbox within a cgroup managed by \code{systemd-run} (\cref{tab:cgroup}). Without the cgroup, the check spawned 500 processes and consumed 360\,MiB in \code{/work}. With the cgroup, only 59 processes were allowed, and the kernel terminated the check via \code{SIGKILL} when it exceeded the memory limit. Note that swap must be disabled for the cgroup, otherwise the kernel may move the files to swap, allowing the write to succeed.

\paragraph{Cost.} An empty check (\code{true}) took 3\,ms without the sandbox and 18\,ms with it (averaged over 150 runs), resulting in a sandbox overhead of approximately 15\,ms. Compiling and running a 144-line Rust program took 360\,ms without the sandbox and 365\,ms with it (20 runs). The overhead is minimal here because the sandbox compiles into memory. When the plain run is also performed in memory (349\,ms), the sandbox again adds approximately 15\,ms, or 4\%.

\paragraph{Hosts.} On an Ubuntu 24.04 host where AppArmor restricts unprivileged user namespaces via \path{kernel.apparmor_restrict_unprivileged_userns=1}, \code{bwrap} failed with \code{loopback: Failed RTM\_NEWADDR: Operation not permitted}. In such cases, the harness terminates before the first request.

\paragraph{A real task.} We requested a Rust program demonstrating \code{std::mem::replace}, \code{std::mem::take}, and other standard library memory functions. The initial attempt failed all 5 retries with \code{EOF while parsing a string}. Every reply reached the model's output limit within the program, and because each retry was as long as the previous one, the error messages provided no useful feedback. A subsequent run succeeded on the first attempt. However, the program compiled with a warning, and a comment incorrectly described \code{std::mem::swap} as atomic. The check failed to verify either the code's correctness or the accuracy of the comments.

\section{Limitations}
\label{sec:limits}

The primary challenge remains the nature of the check itself. Because the model authors the check, a command such as \code{echo ok} can be used to bypass any program-level verification. We intend to replace the arbitrary shell command with a structured, typed description, from which the harness will construct the appropriate command. This would allow the model to specify the intended test while the harness maintains control over the execution parameters.

Three remaining gaps in confinement exist within the current harness. First, check output is returned to the model unmodified, which allows a program to print instructions for subsequent attempts. Second, we do not employ a \emph{seccomp} filter \cite{seccomp2}; consequently, any system call permitted by the user namespace, including those exposing kernel vulnerabilities, remains reachable. Third, the harness does not impose limits on the total number of processes or the aggregate memory consumption. While the cgroup configuration described in \cref{tab:cgroup} mitigates these issues in our tests, it requires \emph{systemd} to delegate the \code{pids} and \code{memory} controllers to the user. Finally, our evaluation was limited to a single hardware platform and a single model.

\section{Related work}

Existing research explores the intersection of language models and tool use. ReAct interleaves reasoning with actions, Toolformer teaches models when to invoke APIs, and CodeAct and SWE-agent enable models to execute code and tests \cite{react2023,toolformer2023,codeact2024,sweagent2024}. These systems focus on the capabilities models gain through tool access. In contrast, our harness reduces the agent loop to a set of fixed steps, ensuring that every effect of a model reply on the host can be explicitly listed and confined. Srinivasan describes runtime patterns for production agents, including shared state machines and supervisor-gate architectures \cite{sdb2026}. Our harness combines these concepts, with the distinction that the model provides the content for the check that the gate executes. MetaAgent utilizes finite state machines to build multi-agent systems with revisitable states \cite{metaagent2025}. Our implementation uses five fixed, hand-written states, ensuring that the transitions are sufficiently few to be verified by inspection.

\section{Conclusion}

We have developed a coding harness designed to trace every execution path from a language model's reply to the host. The agent loop is implemented as a fixed state machine, ensuring the model only controls the content of each step. This content reaches the host as a file with a plain name and as a check executed within a transient namespace sandbox, incurring approximately 15\,ms of overhead per check. Our evaluation demonstrates that no generated escape attempts left a trace on the host, and that cgroups effectively bound the processes and memory that the sandbox alone does not. Future work will focus on removing the model's ability to author the check itself.

\bibliographystyle{plainnat}
\bibliography{refs}

\end{document}
